% ============================================================
%  CHAPTER 4: NEWTON'S LAWS AND DYNAMICS
%  The Physics Codex: A Complete University Foundation
%  Korede Samuel Omotosho (Falcon)
%  Aurikrex Learning Systems, 2026
%
%  File: chapters/part1/ch04_dynamics.tex
% ============================================================

\chapter{Newton's Laws and Dynamics}
\label{ch:dynamics}

\chapterquote{If I have seen further, it is by standing
on the shoulders of giants.}{Isaac Newton}

\begin{objectivebox}
By the end of this chapter, you should be able to:
\begin{enumerate}[noitemsep]
  \item State and explain Newton's three laws of motion
  with real examples
  \item Define linear momentum and impulse
  \item Apply the law of conservation of momentum to
  collisions and explosions
  \item Distinguish between elastic and inelastic
  collisions
  \item Apply Newton's laws to connected bodies, pulleys,
  and inclined planes
  \item Define friction, state its laws, and apply them
  to practical problems
  \item Analyse motion on inclined planes with and
  without friction
  \item Understand and apply the concept of terminal
  velocity
\end{enumerate}
\end{objectivebox}

\vspace{0.4cm}

\minitoc

\vspace{0.5cm}

% ============================================================
\section{From Kinematics to Dynamics}
% ============================================================

Chapter~3 described \textit{how} objects move. This
chapter asks \textit{why}. The answer was given in 1687
by Isaac Newton in his masterwork
\textit{Philosophi\ae\ Naturalis Principia Mathematica}
--- the \textit{Principia} --- widely considered the most
important scientific book ever written. In it, Newton
unified the motion of the planets, the fall of objects
on Earth, the behaviour of tides, and the paths of
comets under a single set of three laws.

\begin{historybox}[Newton and the Apple]
Newton himself claimed that the sight of an apple falling
from a tree in his garden in Woolsthorpe, Lincolnshire,
around 1666, set him thinking about gravity. Whether the
story is embellished or not, what is certain is that
Newton spent roughly twenty years developing his ideas
before finally publishing them. His great insight was not
that gravity existed --- everyone knew things fell --- but
that the same gravitational force that pulled the apple
also kept the Moon in orbit. The Moon, he realised, is
perpetually \textit{falling} --- it just keeps missing
the Earth because of its sideways velocity.
\end{historybox}

% ============================================================
\section{Newton's First Law: Inertia}
% ============================================================

\begin{lawbox}[Newton's First Law of Motion]
A body remains at rest, or continues to move in a
straight line at constant velocity, unless acted upon
by a net (resultant) external force.

Mathematically: if $\sum\vect{F} = \vect{0}$,
then $\vect{v} = \text{constant}$ (including $v = 0$).
\end{lawbox}

This law contains two profound ideas. First, rest and
uniform motion are equivalent states --- a body doing
either requires no explanation. Only \textit{change}
in motion requires explanation. Second, every body has
an intrinsic resistance to change in its state of motion.
This resistance is called \textbf{inertia}.

\begin{definitionbox}[Inertia]
\textbf{Inertia} is the tendency of a body to resist
any change in its state of rest or uniform motion.
The measure of inertia is \textbf{mass}. A more massive
body has more inertia --- it is harder to start, stop,
or change the direction of.
\end{definitionbox}

\begin{realworldbox}[Inertia on a Danfo]
When a Lagos danfo suddenly brakes, passengers who are
standing lurch forward. Their bodies were moving with
the bus and, by Newton's First Law, they tend to
continue moving forward. Only a force (from a handrail,
another passenger, or the seat) stops them. Seatbelts
exist precisely to provide this stopping force before
the windscreen does.
\end{realworldbox}

% ============================================================
\section{Newton's Second Law: Force and Momentum}
% ============================================================

\begin{lawbox}[Newton's Second Law of Motion]
The rate of change of momentum of a body is directly
proportional to the net external force acting on it,
and takes place in the direction of that force.
\[
  \vect{F}_{net} = \frac{d\vect{p}}{dt}
\]
Over a finite interval, the average net force is
\[
  \vect{F}_{\mathrm{avg}} = \frac{\Delta\vect{p}}{\Delta t}
\]
For constant mass:
\[
  \vect{F}_{net} = m\vect{a}
\]
The SI unit of force is the \textbf{newton} (N):
$1\ \text{N} = 1\ \text{kg\,m\,s}^{-2}$
\end{lawbox}

\begin{definitionbox}[Linear Momentum]
The \textbf{linear momentum} $\vect{p}$ of a body is
the product of its mass and velocity:
\[
  \vect{p} = m\vect{v}
\]
Momentum is a \textbf{vector} quantity. SI unit:
kg\,m\,s$^{-1}$ or N\,s.
\end{definitionbox}

\begin{definitionbox}[Impulse]
The \textbf{impulse} $\vect{J}$ of a force is the
product of the force and the time interval over which
it acts:
\[
  \vect{J} = \vect{F}\,\Delta t
  = \Delta\vect{p} = m\vect{v} - m\vect{u}
\]
Impulse equals the change in momentum.
SI unit: N\,s = kg\,m\,s$^{-1}$.
\end{definitionbox}

\begin{keybox}[Impulse-Momentum Theorem]
\[
  \vect{F}\,\Delta t = m\vect{v} - m\vect{u}
\]
A large force acting for a short time produces the same
change in momentum as a small force acting for a long
time. This is why airbags save lives: they increase the
collision time, reducing the peak force on the occupant.
\end{keybox}

\begin{examplebox}[4.1]
\stepcounter{workedexample}
A $1500\ \text{kg}$ car travelling at
$20\ \text{m\,s}^{-1}$ is brought to rest in $4.0\ \text{s}$
by applying the brakes.\\
(a) Find the change in momentum.\\
(b) Find the average braking force.\\
(c) If the same car hit a wall and stopped in $0.1\ \text{s}$,
what would the average force be?
\end{examplebox}

\begin{solutionbox}
\textbf{(a)} Change in momentum:
\[
  \Delta p = mv - mu = 1500(0) - 1500(20)
  = -30\,000\ \text{kg\,m\,s}^{-1}
\]
Magnitude: $30\,000\ \text{N\,s}$

\textbf{(b)} Average braking force:
\[
  F = \frac{\Delta p}{\Delta t}
  = \frac{30\,000}{4.0} = 7\,500\ \text{N}
\]

\textbf{(c)} Hitting the wall ($\Delta t = 0.1\ \text{s}$):
\[
  F = \frac{30\,000}{0.1} = 300\,000\ \text{N}
\]
This is 40 times larger --- enough to be fatal. The
same change in momentum, but 40 times less time, means
40 times the force. This explains why crumple zones,
airbags and crash barriers save lives: they increase
$\Delta t$.
\end{solutionbox}

% ============================================================
\section{Newton's Third Law: Action and Reaction}
% ============================================================

\begin{lawbox}[Newton's Third Law of Motion]
For every action (force), there is an equal and opposite
reaction. If body A exerts a force $\vect{F}$ on body B,
then body B exerts a force $-\vect{F}$ on body A.
\[
  \vect{F}_{AB} = -\vect{F}_{BA}
\]
The two forces are equal in magnitude, opposite in
direction, act on \textbf{different} bodies, are of the
same type, and act simultaneously.
\end{lawbox}

\begin{misconceptionbox}
The action and reaction forces in Newton's Third Law
\textbf{never cancel each other}. They cannot, because
they act on different bodies. If you push a wall with
force $F$, the wall pushes you back with force $F$ ---
but these forces act on different objects (you and the
wall). To find whether something accelerates, add up
all forces on \textit{that object alone}.
\end{misconceptionbox}

\begin{realworldbox}[Rocket Propulsion]
A rocket engine burns fuel, expelling hot gases
\textit{backward} at high speed. By Newton's Third
Law, the gases push the rocket \textit{forward} with
an equal and opposite force (thrust). There is no air
to push against in space --- the rocket pushes against
the expelled gas itself. This is why rockets work in
the vacuum of space, while aeroplane propellers do not.
The NASC (National Space Research and Development
Agency) in Abuja studies rocket propulsion for Nigeria's
space programme.
\end{realworldbox}

% ============================================================
\section{Conservation of Linear Momentum}
% ============================================================

\begin{lawbox}[Law of Conservation of Linear Momentum]
The total linear momentum of a closed system (no net
external force) remains constant, regardless of the
interactions between the bodies within the system.
\[
  \sum \vect{p}_{before}
  = \sum \vect{p}_{after}
\]
\[
  m_1\vect{u}_1 + m_2\vect{u}_2
  = m_1\vect{v}_1 + m_2\vect{v}_2
\]
This law follows directly from Newton's Second and Third
Laws combined.
\end{lawbox}

\subsection{Types of Collision}

\begin{definitionbox}[Elastic and Inelastic Collisions]
\textbf{Elastic collision:} Both momentum AND kinetic
energy are conserved. In practice, perfectly elastic
collisions occur only at the atomic/molecular level.
Billiard balls approximate elastic collisions.

\vspace{0.4cm}

\textbf{Inelastic collision:} Momentum is conserved but
kinetic energy is NOT conserved. Some KE is converted
to heat, sound, and deformation.

\vspace{0.4cm}

\textbf{Perfectly inelastic collision:} The bodies stick
together after impact and move with a common velocity.
Maximum kinetic energy is lost.
\[
  m_1u_1 + m_2u_2 = (m_1+m_2)V
\]
\end{definitionbox}

\begin{examplebox}[4.2]
\stepcounter{workedexample}
A $3.0\ \text{kg}$ trolley moving at
$6.0\ \text{m\,s}^{-1}$ collides and sticks to a
$2.0\ \text{kg}$ stationary trolley.\\
(a) Find the common velocity after collision.\\
(b) Calculate the kinetic energy before and after.\\
(c) What fraction of the kinetic energy is lost?
\end{examplebox}

\begin{solutionbox}
\textbf{(a)} Conservation of momentum (perfectly
inelastic):
\[
  m_1u_1 + m_2u_2 = (m_1+m_2)V
\]
\[
  (3.0)(6.0) + (2.0)(0) = (3.0+2.0)V
\]
\[
  18 = 5V \implies V = 3.6\ \text{m\,s}^{-1}
\]

\textbf{(b)} Kinetic energies:
\[
  KE_{before}
  = \tfrac{1}{2}(3.0)(6.0)^2 = 54\ \text{J}
\]
\[
  KE_{after}
  = \tfrac{1}{2}(5.0)(3.6)^2 = 32.4\ \text{J}
\]

\textbf{(c)} Energy lost $= 54 - 32.4 = 21.6\ \text{J}$
\[
  \text{Fraction lost}
  = \frac{21.6}{54} = 0.40 = 40\%
\]
\end{solutionbox}

\begin{examplebox}[4.3]
\stepcounter{workedexample}
A gun of mass $2.0\ \text{kg}$ fires a bullet of mass
$20\ \text{g}$ at $400\ \text{m\,s}^{-1}$.
Find the recoil velocity of the gun.
\end{examplebox}

\begin{solutionbox}
Initially both gun and bullet are at rest:
total initial momentum $= 0$.

By conservation of momentum:
\[
  0 = m_{bullet}v_{bullet} + m_{gun}v_{gun}
\]
\[
  0 = (0.020)(400) + (2.0)v_{gun}
\]
\[
  v_{gun} = -\frac{8.0}{2.0} = -4.0\ \text{m\,s}^{-1}
\]
The gun recoils at $4.0\ \text{m\,s}^{-1}$ in the
direction opposite to the bullet. The negative sign
confirms this.
\end{solutionbox}

% ============================================================
\section{Applications of Newton's Laws}
% ============================================================

\subsection{Connected Bodies and Pulleys}

\begin{examplebox}[4.4]
\stepcounter{workedexample}
Two blocks, $A$ ($4.0\ \text{kg}$) and $B$
($6.0\ \text{kg}$), are connected by a light,
inextensible string over a smooth pulley (Atwood's
machine). Find: (a) the acceleration; (b) the tension.
($g = 10\ \text{m\,s}^{-2}$)
\end{examplebox}

\begin{solutionbox}
\begin{center}
\begin{tikzpicture}[scale=0.8]
  % Pulley
  \draw[thick] (0,2) circle (0.5cm);
  \fill[gray!30] (0,2) circle (0.5cm);
  \draw[thick] (-0.5,2.5) -- (0.5,2.5);

  % String
  \draw[thick] (-0.5,2) -- (-0.5,0.5);
  \draw[thick] (0.5,2) -- (0.5,0.5);

  % Block A (left, lighter)
  \fill[codexblue!30] (-1,0) rectangle (0,0.5);
  \node at (-0.5,0.25) {$A$};
  \node[left, font=\small] at (-1,0.25) {$4$ kg};
  \draw[->, thick, codexblue] (-0.5,0)
    -- (-0.5,-0.8) node[left]{$m_Ag$};
  \draw[->, thick, codexred] (-0.5,0.5)
    -- (-0.5,1.2) node[left]{$T$};

  % Block B (right, heavier)
  \fill[codexred!30] (0,0) rectangle (1,0.5);
  \node at (0.5,0.25) {$B$};
  \node[right, font=\small] at (1,0.25) {$6$ kg};
  \draw[->, thick, codexblue] (0.5,0)
    -- (0.5,-0.8) node[right]{$m_Bg$};
  \draw[->, thick, codexred] (0.5,0.5)
    -- (0.5,1.2) node[right]{$T$};

  % Acceleration arrows
  \draw[->, thick, codexgreen, dashed]
    (-0.5,-1.0) -- (-0.5,-0.5)
    node[left, font=\small]{$a$ (up)};
  \draw[->, thick, codexgreen, dashed]
    (0.5,-1.0) -- (0.5,-1.5)
    node[right, font=\small]{$a$ (down)};
\end{tikzpicture}
\end{center}

$B$ is heavier so it descends and $A$ ascends.
Take downward for $B$ as positive, upward for $A$
as positive.

For block $B$ (net downward force):
\[
  m_B g - T = m_B a \quad\cdots(1)
\]
For block $A$ (net upward force):
\[
  T - m_A g = m_A a \quad\cdots(2)
\]

Adding equations (1) and (2):
\[
  (m_B - m_A)g = (m_A + m_B)a
\]
\[
  a = \frac{(m_B - m_A)g}{m_A + m_B}
  = \frac{(6-4)(10)}{4+6}
  = \frac{20}{10} = 2.0\ \text{m\,s}^{-2}
\]

Substituting into (2):
\[
  T = m_A(g + a) = 4.0(10 + 2.0) = 48\ \text{N}
\]
\end{solutionbox}

\subsection{Motion on an Inclined Plane}

\begin{diagbox}[Forces on an Inclined Plane]
\begin{center}
\begin{tikzpicture}[scale=1.0]
  % Incline
  \fill[gray!20] (0,0) -- (5,0) -- (5,3) -- cycle;
  \draw[thick] (0,0) -- (5,0) -- (5,3) -- cycle;
  \draw[codexgold] (0.6,0) arc(0:31:0.6);
  \node[codexgold] at (1.0,0.2) {$\theta$};

  % Block
  \begin{scope}[rotate=31, shift={(2.5,0.15)}]
    \fill[codexblue!30] (0,0) rectangle (1,0.7);
    \node at (0.5,0.35) {$m$};
  \end{scope}

  % Forces from centre of block
  \coordinate (C) at (3.0, 1.5);

  % Weight
  \draw[->, very thick, codexred]
    (C) -- ++(0,-1.8)
    node[below]{$W = mg$};

  % Normal reaction (perpendicular to surface)
  \draw[->, very thick, codexblue]
    (C) -- ++(-0.95, 1.56)
    node[above left]{$R$};

  % Component parallel to incline (down the slope)
  \draw[->, thick, codexgreen, dashed]
    (C) -- ++(0.95,-0.57)
    node[right, font=\small]{$mg\sin\theta$};

  % Component perpendicular to incline
  \draw[->, thick, codexgold, dashed]
    (C) -- ++(-0.57,-0.95)
    node[left, font=\small]{$mg\cos\theta$};
\end{tikzpicture}
\end{center}
\end{diagbox}

For a body on a frictionless incline at angle $\theta$:
\[
  R = mg\cos\theta \quad \text{(perpendicular to surface)}
\]
\[
  F_{net} = mg\sin\theta \quad \text{(along surface, downward)}
\]
\[
  a = g\sin\theta \quad \text{(acceleration down the incline)}
\]

With friction (coefficient $\mu$):
\[
  F_{friction} = \mu R = \mu mg\cos\theta
\]
\[
  a = g(\sin\theta - \mu\cos\theta)
  \quad \text{(sliding down, friction acts up)}
\]

\begin{examplebox}[4.5]
\stepcounter{workedexample}
A $5.0\ \text{kg}$ crate slides down a ramp inclined
at $30°$ to the horizontal. The coefficient of kinetic
friction between the crate and the ramp is $0.20$.
Find: (a) the normal reaction; (b) the friction force;
(c) the acceleration.
($g = 10\ \text{m\,s}^{-2}$)
\end{examplebox}

\begin{solutionbox}
\textbf{(a)} Normal reaction:
\[
  R = mg\cos\theta = 5.0 \times 10 \times \cos 30°
  = 50 \times 0.866 = 43.3\ \text{N}
\]

\textbf{(b)} Friction force:
\[
  f = \mu R = 0.20 \times 43.3 = 8.66\ \text{N}
\]
(acting up the incline, opposing motion)

\textbf{(c)} Net force down the incline:
\[
  F_{net} = mg\sin\theta - f
  = 50\sin 30° - 8.66
  = 25 - 8.66 = 16.3\ \text{N}
\]
\[
  a = \frac{F_{net}}{m} = \frac{16.3}{5.0}
  = 3.3\ \text{m\,s}^{-2}
\]
\end{solutionbox}

% ============================================================
\section{Friction}
% ============================================================

\begin{definitionbox}[Friction]
\textbf{Friction} is the force that opposes the relative
motion (or tendency of motion) between two surfaces in
contact. It acts tangentially to the surfaces, opposing
the direction of motion or impending motion.
\end{definitionbox}

\begin{lawbox}[Laws of Solid Friction (Amontons-Coulomb)]
\begin{enumerate}[noitemsep]
  \item Friction force is proportional to the normal
  reaction: $F \propto R$, i.e.\ $F = \mu R$
  \item Friction is independent of the area of contact
  between the surfaces (for given $R$)
  \item Kinetic friction is independent of the relative
  speed of the surfaces
  \item Static friction is always greater than or equal
  to kinetic friction:
  $\mu_s \geq \mu_k$
\end{enumerate}
where $\mu$ is the \textbf{coefficient of friction}
(dimensionless).
\end{lawbox}

\begin{definitionbox}[Static and Kinetic Friction]
\textbf{Static friction} $f_s$ prevents two surfaces
from sliding relative to each other. It has a maximum
value: $f_{s,max} = \mu_s R$.

\textbf{Kinetic (dynamic) friction} $f_k$ acts between
surfaces already in relative motion: $f_k = \mu_k R$.

Always: $f_{s,max} \geq f_k$, so $\mu_s \geq \mu_k$.
\end{definitionbox}

% ============================================================
\section{Terminal Velocity}
% ============================================================

When an object falls through a fluid (air, water), the
fluid exerts a \textbf{drag force} that opposes the
motion and increases with speed. At some point, the drag
force equals the weight. At that point, the net force
is zero and the object stops accelerating --- it falls
at constant speed called the \textbf{terminal velocity}.

\begin{processbox}[Stages of Free Fall with Air Resistance]
\begin{enumerate}[noitemsep]
  \item \textbf{Initial:} Object starts from rest.
  Weight $> $ drag. Net downward force is large.
  Acceleration is close to $g$.
  \item \textbf{Intermediate:} As speed increases,
  drag increases. Net downward force decreases.
  Acceleration decreases but body still speeds up.
  \item \textbf{Terminal:} Speed reaches terminal
  velocity $v_T$ where drag $=$ weight.
  Net force $= 0$. Acceleration $= 0$.
  Speed is constant.
\end{enumerate}
\[
  \text{At terminal velocity: } F_{drag} = mg
\]
\end{processbox}

\begin{diagbox}[Velocity-Time Graph for Terminal Velocity]
\begin{center}
\begin{tikzpicture}
\begin{axis}[
  width=10cm, height=6cm,
  xlabel={Time $t$ (s)},
  ylabel={Velocity $v$ (m\,s$^{-1}$)},
  xmin=0, xmax=10, ymin=0, ymax=60,
  ytick={0,10,20,30,40,50,55},
  grid=major, grid style={gray!20},
  yticklabels={0,10,20,30,40,50,$v_T$},
]
  % Terminal velocity curve (approaching asymptote)
  \addplot[codexblue, very thick, domain=0:10, samples=100]
    {55*(1 - exp(-0.4*x))};

  % Terminal velocity line
  \addplot[codexred, thick, dashed, domain=0:10]{55};
  \node[codexred, right] at (axis cs:9.5,55)
    {$v_T$};

  % Annotation: steep initial gradient
  \draw[->, gray] (axis cs:1,20) -- (axis cs:0.5,10)
    node[left, font=\tiny]{High initial $a$};

  % Annotation: gradient decreasing
  \draw[->, gray] (axis cs:4,45) -- (axis cs:5,50)
    node[right, font=\tiny]{$a$ decreasing};

  % Annotation: flat at terminal
  \draw[->, gray] (axis cs:8,53) -- (axis cs:9,55)
    node[right, font=\tiny]{$a = 0$};
\end{axis}
\end{tikzpicture}
\end{center}
\end{diagbox}

\begin{realworldbox}[Parachutes and Terminal Velocity]
A skydiver in free fall (no parachute) reaches a
terminal velocity of approximately $55\ \text{m\,s}^{-1}$
($200\ \text{km\,h}^{-1}$). When the parachute opens,
it enormously increases the drag force, decelerating
the skydiver rapidly to a new, much lower terminal
velocity of about $5\ \text{m\,s}^{-1}$ --- safe for
landing. The parachute works not by eliminating gravity
but by increasing drag until it equals the weight at a
safe speed.
\end{realworldbox}

% ============================================================
\section{Weight, Mass and the Normal Force}
% ============================================================

\begin{definitionbox}[Weight]
The \textbf{weight} $W$ of a body is the gravitational
force exerted on it by the Earth (or whatever
gravitational body it is near):
\[
  W = mg
\]
Weight is a \textbf{vector} (directed toward the centre
of the Earth). SI unit: newton (N).

Weight depends on location. The same $5.0\ \text{kg}$
mass weighs $49\ \text{N}$ on Earth, $8.2\ \text{N}$
on the Moon, and $0\ \text{N}$ in deep space.
Mass is constant everywhere.
\end{definitionbox}

\begin{definitionbox}[Normal Force]
The \textbf{normal force} (or normal reaction) $N$ is
the contact force exerted by a surface on an object,
perpendicular to the surface. It prevents objects from
passing through surfaces. It is not always equal to
weight --- it depends on the situation.
\end{definitionbox}

% ============================================================
\section{Chapter Summary}
% ============================================================

\begin{summarybox}
\textbf{Newton's First Law:} Body at rest or constant
velocity unless net force acts. Defines inertia.

\textbf{Newton's Second Law:} $\vect{F}_{net} = m\vect{a}$
or $\vect{F} = \Delta\vect{p}/\Delta t$.

\textbf{Newton's Third Law:} Action = $-$Reaction
(on different bodies).

\textbf{Momentum:} $\vect{p} = m\vect{v}$ (vector, N\,s)

\textbf{Impulse:} $\vect{J} = \vect{F}\Delta t
= \Delta\vect{p}$

\textbf{Conservation of momentum:}
$\sum\vect{p}_{before} = \sum\vect{p}_{after}$
(closed system)

\textbf{Elastic collision:} momentum and KE conserved.
\textbf{Perfectly inelastic:}
$(m_1+m_2)V = m_1u_1+m_2u_2$

\textbf{Inclined plane:}
$N = mg\cos\theta$;
$a = g(\sin\theta - \mu\cos\theta)$

\textbf{Friction:} $f = \mu N$; static $\geq$ kinetic.

\textbf{Terminal velocity:} drag = weight; $a = 0$.

\textbf{Weight:} $W = mg$ (vector, varies with location).
\textbf{Mass:} constant everywhere.
\end{summarybox}

% ============================================================
\newpage
\begin{exercisebox}[4]

\subsection*{Section A: Multiple Choice}

\textbf{A1.} A force of $20\ \text{N}$ acts on a
$4.0\ \text{kg}$ mass for $5.0\ \text{s}$.
The change in momentum is:

\mcoption{A}{$4\ \text{N\,s}$}
\mcoption{B}{$16\ \text{N\,s}$}
\mcoption{C}{$100\ \text{N\,s}$}
\mcoption{D}{$400\ \text{N\,s}$}
\ansref{Ch4-A1}

\vspace{0.4cm}

\textbf{A2.} A $2\ \text{kg}$ and a $6\ \text{kg}$ mass
hang over a smooth pulley. The acceleration of the
system is (take $g = 10\ \text{m\,s}^{-2}$):

\mcoption{A}{$2.5\ \text{m\,s}^{-2}$}
\mcoption{B}{$5.0\ \text{m\,s}^{-2}$}
\mcoption{C}{$7.5\ \text{m\,s}^{-2}$}
\mcoption{D}{$10\ \text{m\,s}^{-2}$}
\ansref{Ch4-A2}

\vspace{0.4cm}

\textbf{A3.} Newton's Third Law pairs are:

\mcoption{A}{Weight of a book and the normal force
from the table on the book}
\mcoption{B}{Weight of a book and the gravitational
pull of the book on the Earth}
\mcoption{C}{Tension in a rope and the weight of
the hanging object}
\mcoption{D}{Friction and normal force on a
sliding block}
\ansref{Ch4-A3}

\vspace{0.4cm}

\textbf{A4.} A $5\ \text{kg}$ block slides down a
smooth incline of $37°$. Its acceleration is
(sin$37° = 0.6$, $g = 10\ \text{m\,s}^{-2}$):

\mcoption{A}{$4\ \text{m\,s}^{-2}$}
\mcoption{B}{$6\ \text{m\,s}^{-2}$}
\mcoption{C}{$8\ \text{m\,s}^{-2}$}
\mcoption{D}{$10\ \text{m\,s}^{-2}$}
\ansref{Ch4-A4}

\vspace{0.4cm}

\textbf{A5.} An object reaches terminal velocity when:

\mcoption{A}{Its weight is zero}
\mcoption{B}{Its acceleration equals $g$}
\mcoption{C}{The drag force equals its weight}
\mcoption{D}{Its kinetic energy is maximum}
\ansref{Ch4-A5}

\vspace{0.4cm}

\textbf{A6.} In a perfectly inelastic collision:

\mcoption{A}{Both momentum and KE are conserved}
\mcoption{B}{Momentum is conserved but KE is not}
\mcoption{C}{KE is conserved but momentum is not}
\mcoption{D}{Neither momentum nor KE is conserved}
\ansref{Ch4-A6}

\vspace{0.4cm}

\textbf{A7.} The coefficient of friction between
two surfaces is $0.4$. A $10\ \text{kg}$ block rests
on a horizontal surface. The friction force when a
horizontal force of $30\ \text{N}$ is applied
(take $g = 10\ \text{m\,s}^{-2}$) is:

\mcoption{A}{$4\ \text{N}$}
\mcoption{B}{$30\ \text{N}$}
\mcoption{C}{$40\ \text{N}$}
\mcoption{D}{$12\ \text{N}$}
\ansref{Ch4-A7}

\vspace{0.4cm}

\textbf{A8.} A $60\ \text{kg}$ person stands on a
scale in a lift accelerating upward at
$2\ \text{m\,s}^{-2}$. The scale reads:
($g = 10\ \text{m\,s}^{-2}$)

\mcoption{A}{$480\ \text{N}$}
\mcoption{B}{$600\ \text{N}$}
\mcoption{C}{$720\ \text{N}$}
\mcoption{D}{$120\ \text{N}$}
\ansref{Ch4-A8}

\vspace{0.4cm}

\textbf{A9.} A rifle of mass $3.0\ \text{kg}$ fires
a $15\ \text{g}$ bullet at $600\ \text{m\,s}^{-1}$.
The recoil speed of the rifle is:

\mcoption{A}{$0.5\ \text{m\,s}^{-1}$}
\mcoption{B}{$1.5\ \text{m\,s}^{-1}$}
\mcoption{C}{$3.0\ \text{m\,s}^{-1}$}
\mcoption{D}{$6.0\ \text{m\,s}^{-1}$}
\ansref{Ch4-A9}

\vspace{0.4cm}

\textbf{A10.} Which of the following is a consequence
of Newton's First Law?

\mcoption{A}{A falling body accelerates at $g$}
\mcoption{B}{A book on a table is at rest because
the table exerts a normal force}
\mcoption{C}{Passengers lurch forward when a bus
brakes suddenly}
\mcoption{D}{A rocket accelerates when fuel is burned}
\ansref{Ch4-A10}

\vspace{0.6cm}

\subsection*{Section B: Theory and Calculations}

\textbf{B1.} A $1200\ \text{kg}$ car travelling at
$30\ \text{m\,s}^{-1}$ rear-ends a stationary
$800\ \text{kg}$ car. After the collision, the two
cars move together.\\[4pt]
(a) Find the common velocity after collision.\\
(b) Calculate the kinetic energy lost in the collision.\\
(c) Express this as a percentage of the original KE.\\
(d) In what form does this lost energy appear?
\hfill\ansref{Ch4-B1}

\vspace{0.4cm}

\textbf{B2.} A $8.0\ \text{kg}$ block is placed on an
incline of angle $25°$. The coefficient of static
friction is $0.45$ and kinetic friction is $0.35$.
($g = 10\ \text{m\,s}^{-2}$, sin$25° = 0.423$,
cos$25° = 0.906$)\\[4pt]
(a) Does the block slide? Justify using force
comparison.\\
(b) If a force of $20\ \text{N}$ acts down the incline,
find the net acceleration.\\
(c) If the block is already sliding, find its
acceleration without the extra force.
\hfill\ansref{Ch4-B2}

\vspace{0.4cm}

\textbf{B3.} Two masses, $m_1 = 3.0\ \text{kg}$ on a
horizontal frictionless table and $m_2 = 5.0\ \text{kg}$
hanging over the edge, are connected by a light string
over a smooth pulley at the edge of the table.\\[4pt]
(a) Draw a free-body diagram for each mass.\\
(b) Write Newton's second law equations for each mass.\\
(c) Solve for the acceleration of the system.\\
(d) Find the tension in the string.\\
(e) If friction exists with $\mu_k = 0.20$, what is
the new acceleration?
\hfill\ansref{Ch4-B3}

\vspace{0.4cm}

\textbf{B4.} A parachutist of total mass $90\ \text{kg}$
(including equipment) falls from rest. The drag force
at velocity $v$ is given by $F_d = kv^2$ where
$k = 0.6\ \text{kg\,m}^{-1}$. ($g = 10\ \text{m\,s}^{-2}$)\\[4pt]
(a) Calculate the terminal velocity.\\
(b) At $v = 20\ \text{m\,s}^{-1}$, calculate the net
force and the acceleration.\\
(c) Sketch the velocity-time graph from release to
terminal velocity, labelling key features.\\
(d) Explain what happens to the acceleration as the
speed increases from zero to terminal velocity.
\hfill\ansref{Ch4-B4}

\vspace{0.4cm}

\textbf{B5.} A $500\ \text{kg}$ cannon fires a
$10\ \text{kg}$ shell horizontally. The shell leaves
the barrel at $400\ \text{m\,s}^{-1}$. The cannon
recoils on a frictionless surface.\\[4pt]
(a) Find the recoil velocity of the cannon.\\
(b) After firing, the cannon rolls back $2.0\ \text{m}$
before a brake stops it. What is the average braking
force?\\
(c) How does the impulse on the shell compare with
the impulse on the cannon? Explain using Newton's
Third Law.\\
(d) The cannon and shell experience the same impulse
magnitude but very different velocities. Why?
\hfill\ansref{Ch4-B5}

\end{exercisebox}
